\section*{General Conclusion}

This thesis presented the design and implementation of a complete ride-sharing platform, developed as a backend-driven, real-time, and intelligent system. The objective was not only to build a functional application, but to address the limitations of traditional ride-sharing approaches by introducing flexibility, synchronization, and predictive capabilities within a unified architecture.

The system was constructed progressively through an iterative development process, where each phase contributed to a distinct layer of functionality. The initial stages established a reliable foundation through controlled user management and a structured ride lifecycle, ensuring that all operations remain consistent and validated at the backend level. This foundation was then extended with a real-time communication layer based on an event-driven model, enabling continuous synchronization between users and transforming the platform into a dynamic, coordinated environment.

Building on this architecture, the system introduced intelligent components that enhance both adaptability and usability. The smart matching mechanism reformulated the ride allocation problem as a trajectory-based optimization task, allowing flexible pickup points and significantly increasing matching opportunities. In parallel, the machine learning component added a predictive dimension by identifying recurring commuting patterns and generating context-aware suggestions. The deliberate separation between deterministic logic and predictive models ensures that system reliability is preserved while enabling intelligent assistance.

From an architectural perspective, the platform reflects a set of deliberate design choices. Centralizing business logic in the backend guarantees consistency, security, and control over system behavior. The integration of real-time communication through WebSocket technology enables efficient state propagation without reliance on polling mechanisms. Additionally, the inclusion of an observability layer provides transparency into system operations, supporting monitoring, debugging, and overall system reliability.

The final evaluation confirmed that the system operates coherently as an integrated solution. End-to-end workflows are executed reliably, real-time interactions remain consistent across multiple clients, and intelligent components provide meaningful support to users. These results demonstrate that the platform successfully combines correctness, responsiveness, and adaptability within a single system.

Despite these achievements, several limitations remain. The machine learning component relies on synthetic data, which may not fully capture the variability of real-world user behavior. Furthermore, large-scale performance and scalability aspects were not explored in depth, particularly in distributed or high-load environments. Network conditions may also impact the effectiveness of real-time synchronization in less stable contexts.

Future work can extend this system in multiple directions. Integrating real-world datasets would significantly enhance the accuracy and relevance of predictive models. The architecture could be evolved toward a distributed or microservices-based design to support scalability and fault tolerance. Additional intelligent features, such as dynamic pricing, user preference learning, or advanced route optimization, could further improve system effectiveness. Finally, deployment in a real-world environment would provide valuable insights into user behavior and system performance under practical conditions.

In conclusion, this work demonstrates that combining a backend-driven architecture, real-time synchronization, and hybrid intelligent mechanisms enables the development of a robust and adaptive ride-sharing platform. The system not only fulfills its functional requirements but also establishes a solid foundation for future enhancements and real-world applicability.
